\documentclass{article}

\begin{document}

A matemática está presente em várias situações do dia a dia, desde calcular o troco até estimar distâncias e medidas.  
Ela nos ajuda a entender padrões e relações entre quantidades.
Podemos representar operações simples utilizando símbolos matemáticos. Por exemplo, a soma de dois números a e b pode ser escrita como a + b, e uma fração  como a + c/b.
Outra forma de mostrar operações é usando o ambiente de equação para destacá-las:

x1 + y1 = 10
2 / z = 8

Além disso, podemos representar frações, potências e raízes de maneira clara:

a/b, x2, raiz(y)
A matemática não serve apenas para cálculos abstratos; ela aparece em situações concretas, como:

    Calcular a  média das notas escolares.  
    Medir a quantidade de ingredientes em uma receita.  
    Estimar a distância entre dois pontos em um mapa.

Outras aplicações importantes incluem:

    Determinar uma sequência usando u(n+1)=2u(n)+3.  
    Calcular juros simples com J = P x i x n.(Na fórmula de juros simples, P representa o principal, i a taxa e n o número de períodos)
    Resolver problemas de geometria, como encontrar a área da hipotenusa c = raiz(a2 + b2). 


\end{document}
